\section*{Chapter \ref{ch:m3:a-map-of-all-markets} --- A Map of All Markets}

\begin{solution}{exo:m3:a-map-of-all-markets:1}
$1{,}203.6/31{,}800 = 3.8\%$ of the outstanding stock a day.
\end{solution}

\begin{solution}{exo:m3:a-map-of-all-markets:2}
New York is four hours behind UTC in the northern summer: the core session runs 13:30--20:00 UTC and the reference-rate window 19:00--20:00 UTC.
\end{solution}

\begin{solution}{exo:m3:a-map-of-all-markets:3}
For example: E-mini futures at the order-book pole (central limit order book, central counterparty); US corporate bonds towards the dealer pole
(request for quote, bilateral settlement, with growing all-to-all platforms); an \omterm{def:m3:automated-market-makers:amm}{automated market maker} as a third kind, a protocol with neither a
dealer nor a clearing house.
\end{solution}

\begin{solution}{exo:m3:a-map-of-all-markets:4}
$\lceil 24/7 \rceil = 4$ shifts of eight hours; $\lceil 24/11 \rceil = 3$ of twelve.
\end{solution}

\begin{solution}{exo:m3:a-map-of-all-markets:5}
The measures differ: the BIS counts all FX instruments, net of double counting, in one month; the crypto figure is one day on one venue, in one
instrument, reported by the venue itself; and turnover in derivatives is notional, not value exchanged.
\end{solution}

\begin{solution}{exo:m3:a-map-of-all-markets:6}
Its positions are held briefly: most volume comes from market makers and short-horizon traders turning over inventory, not from holders of views
for weeks.
\end{solution}

\begin{solution}{exo:m3:a-map-of-all-markets:7}
Weekday: shifts 00--09, 08--17 and 16--01 UTC with 2, 3 and 3 traders, 8 trader-shifts; Sunday: the same shifts with one trader each, 3.
\end{solution}

\begin{solution}{exo:m3:a-map-of-all-markets:8}
Size is not profit: the largest markets are the most competitive and served by the largest balance sheets; opportunity lies where a firm's skills
meet participants who trade for reasons other than price, at a scale it can serve.
\end{solution}

\begin{solution}{pb:m3:a-map-of-all-markets:1}
\textbf{1.} Three shifts of nine hours overlapping by one. \textbf{2.} 00:00--09:00, 08:00--17:00 and 16:00--01:00 UTC. \textbf{3.} 2, 3 and 3
(crypto plus power at night from 05:00; crypto, power and US in the day; crypto and US in the evening). \textbf{4.} One per shift. \textbf{5.} $5
\times 8 + 2 \times 3 = 46$. \textbf{6.} First shift: funding at 00:00 and at 08:00 with the options expiry (in the overlap); second: 08:00, the
day-ahead auction at 10:00, the LME's official prices and the US open, funding at 16:00 (in the overlap); third: 16:00, the reference-rate window
and the US close at 20:00, and funding at midnight (in the overlap). \textbf{7.} The second. \textbf{8.} The third. \textbf{9.} Open positions and
orders, limits used, pending transfers and margin calls, incidents, and the fixed moments of the next shift. \textbf{10.} In winter the US session
and the reference-rate window move an hour later in UTC and the European auction an hour later, so the second and third shifts' contents shift and
the desk windows should be restated in local time. \textbf{11.} Two shifts would leave two hours uncovered with a one-hour overlap
($\lceil 24/11 \rceil = 3$): twelve-hour shifts still need three, at the cost of fatigue. \textbf{12.} Yes, for the night shift's crypto and early
power coverage, at the cost of a second office and cross-time-zone handovers. \textbf{13.} Automated quoting within limits, hedging and monitoring;
humans must watch risk limits, incidents, venue problems and large moves. \textbf{14.} A reduced crypto rota with automation and a human on call.
\textbf{15.} 46 trader-shifts a week, about nine and a half full-time traders at five shifts each, times the stated cost, plus the overlap and
cover for leave. \textbf{16.} The one whose profit per trader-hour is lowest, measured, not guessed; crypto costs most to staff, the US desk least.
\textbf{17.} Towards markets that reuse the same technology and hours (for instance crypto derivatives on regulated venues) or fill the idle hours.
\textbf{18.} High-velocity markets reward the firm's speed and inventory turnover; low-velocity ones need balance sheet it may not have.
\textbf{19.} \emph{Named result:} three nine-hour shifts a day (00--09, 08--17, 16--01 UTC) with 2, 3 and 3 traders on weekdays and one at weekends.
\textbf{20.} To choose where to trade, and to see what each market demands before going there.
\end{solution}

\begin{solution}{iq:m3:a-map-of-all-markets:1}
FX first, then Treasuries, then equities and crypto; the ranking hides different measures (notional or value, net or gross, one venue or all),
velocity, the share of that turnover a firm can access, and the margins earned.
\iqlookfor{measures, velocity and accessibility.}
\end{solution}

\begin{solution}{iq:m3:a-map-of-all-markets:2}
Positions or orders not understood by the next desk, limits inconsistent across time zones, events in the gap between shifts, differences in local
practice, and fatigue; handover checklists, one book of record and overlap reduce them.
\iqlookfor{information loss at handover.}
\end{solution}

\begin{solution}{iq:m3:a-map-of-all-markets:3}
Volume the firm could trade times the spread it could capture net of adverse selection and fees, from the market's depth, participants and competition,
checked against the profits of firms already there, and discounted for the cost of access.
\iqlookfor{volume times net capture, and competition.}
\end{solution}

\begin{solution}{iq:m3:a-map-of-all-markets:4}
One where participants trade for reasons other than price, technology matters more than balance sheet, and access is cheap: for example market making
in a young derivatives market, with the caveats of venue risk.
\iqlookfor{a reasoned match of edge and constraints.}
\end{solution}

\begin{solution}{iq:m3:a-map-of-all-markets:5}
A calendar service holding each market's session rules in local time, holidays and early closes, converting to UTC for any date with time-zone data, exposing
``open now'' and ``next event'' queries, versioned and tested against the markets' published calendars.
\iqlookfor{local rules, time-zone conversion and versioning.}
\end{solution}

\begin{solution}{iq:m3:a-map-of-all-markets:6}
Choose markets that share infrastructure (order-book market making), fill complementary hours, and have low correlation of profit: for example equity index
futures, crypto perpetuals and European power intraday; justify by edge, hours and capital.
\iqlookfor{shared technology, complementary hours, diversified profits.}
\end{solution}
